\begin{table}[htbp]\centering\begin{singlespace}\small\setlength{\tabcolsep}{3pt}\renewcommand{\arraystretch}{1.3}
\caption{Absolute attention-contact correspondence and paired controls}
\begin{tabular}{>{\raggedright\arraybackslash}p{0.391304347826087\linewidth}>{\raggedright\arraybackslash}p{0.11739130434782608\linewidth}>{\raggedright\arraybackslash}p{0.29347826086956524\linewidth}>{\raggedright\arraybackslash}p{0.09782608695652174\linewidth}}
\toprule
\textbf{Metric} & \textbf{Mean} & \textbf{95\% target-group interval} & \textbf{Groups}\\\midrule
Residue attention AP & 0.13114 & [0.10251, 0.16419] & 10\\
Contact prevalence & 0.07148 & [0.06072, 0.08214] & 10\\
Permutation AP & 0.09311 & [0.07930, 0.10823] & 10\\
Gradient-times-input AP & 0.11911 & [0.09900, 0.14283] & 10\\
Attention minus prevalence & 0.05967 & [0.03084, 0.09449] & 10\\
Attention minus permutation & 0.03803 & [0.00655, 0.07484] & 10\\
Attention minus gradient & 0.01203 & [-0.00614, 0.03359] & 10\\
\bottomrule\end{tabular}\par\vspace{6pt}
{\raggedright\footnotesize Thirty complexes, ten equally weighted target groups, 4-angstrom heavy-atom contacts, 1000 target-group bootstrap resamples. Permutation AP is the per-complex mean over shuffled residue scores. Gradient-times-input scores traverse all protein branches. Paired control increments are secondary analyses of frozen maps. AP and prevalence are unitless; overlapping intervals are not a test of paired differences.\par}
\end{singlespace}\end{table}
